% ============================================================================
\chapter{Les bases de JavaScript}
\label{ch:js-bases}
% ============================================================================

\begin{objectif}
\begin{itemize}
  \item Ex\'ecuter du JavaScript dans un navigateur, dans la console, et dans Node.js.
  \item Ma\^itriser variables, types, op\'erateurs, conversion de types et tests de v\'erit\'e.
  \item \'Ecrire des conditions, des boucles et des fonctions dans toutes les formes modernes.
  \item Comprendre la port\'ee (\emph{scope}), le hoisting et les fermetures (\emph{closures}).
  \item Reconna\^itre les pi\`eges de typage faible (\verb|==|, \verb|NaN|, \verb|this|).
\end{itemize}
\end{objectif}

\section{O\`u et comment ex\'ecuter du JavaScript}

JavaScript est un \emph{moteur}~: V8 dans Chrome et Node, SpiderMonkey dans Firefox,
JavaScriptCore dans Safari. Le langage est d\'efini par la norme \textbf{ECMAScript}
(\link{ECMA-262}{https://262.ecma-international.org/}), \'edit\'ee chaque ann\'ee~; les fonctionnalit\'es
d'ES2015 (dite ES6) \`a ES2024 sont aujourd'hui disponibles partout dans les navigateurs r\'ecents.

\begin{tabz}{llX}
  \headrow \thd{Terrain de jeu} \& \thd{Comment} \& \thd{Id\'eal pour} \\
  \midrule
  Console du navigateur \& \keys{F12} ou \keys{Ctrl+Shift+J} \& Explorer, tester une expression, apprendre le langage. \\
  Fichier HTML \& \verb|<script src="app.js"></script>| \& Tout ce qui touche la page (chapitres DOM). \\
  \link{Playground en ligne}{https://playcode.io/} / \link{CodePen}{https://codepen.io/} / \link{JSFiddle}{https://jsfiddle.net/} \& Ouvrir l'URL \& Partager un bug, prototyper. \\
  Node.js \& \cmd{node fichier.js} ou \cmd{node} (REPL) \& Logique, scripts, outillage, serveurs. \\
  \cmd{deno run} / \cmd{bun} \& \link{deno.com}{https://deno.com/}, \link{bun.sh}{https://bun.sh/} \& Alternatives modernes, rapide, sans configuration. \\
\end{tabz}

\begin{terminal}
# Le REPL (Read-Eval-Print Loop) : la calculatrice geante
$ node
> 2 ** 10
1024
> "nour".toUpperCase()
'NOUR'
> .exit

# Un fichier
$ echo 'console.log("Bonjour le monde");' > hello.js
$ node hello.js
Bonjour le monde
\end{terminal}

\section{Sortir et entrer de l'information}

\begin{codejs}
// En page : la sortie est la console du navigateur
console.log("normal");            // trace
console.info("info"); console.debug("detaill\'e");
console.warn("attention"); console.error("erreur");
console.table([{ a: 1, b: 2 }, { a: 3, b: 4 }]);   // un tableau joliment affiche
console.group("Etapes"); console.log("1"); console.groupEnd();
console.count("appelee");                            // compte les appels
console.time("boucle"); for (let i = 0; i < 1e6; i++); console.timeEnd("boucle");

// Dans un terminal, avec Node :
process.stdout.write("sans saut de ligne");
const reponse = await import("node:readline/promises")
  .then((m) => m.createInterface({ input: process.stdin, output: process.stdout }))
  .then((rl) => rl.question("Ton pr\'enom ? "));
console.log(`Salut ${reponse}`);
\end{codejs}

\section{Variables : \texttt{let}, \texttt{const} (et pourquoi pas \texttt{var})}

\begin{codejs}
const MAX = 100;          // lien immuable : on ne peut pas re-affectuer
let compteur = 0;         // peut changer
var ancienStyle = true;   // ne l'utilise plus : portee de fonction, hoisting piegeux

// "const" ne rend pas la VALEUR immutable, seulement la RE-ASSIGNATION impossible :
const user = { name: "Nour" };
user.name = "Awa";          // parfaitement legal
user = {};                  // TypeError: Assignment to constant variable.
\end{codejs}

R\`egle d'\'ecriture~: \texttt{const} par d\'efaut, \texttt{let} quand la r\'eassignation est n\'ecessaire,
\texttt{var} jamais. Nomme en \verb|camelCase| les variables et fonctions, en
\verb|SCREAMING\_SNAKE\_CASE| les constantes de module, en \verb|PascalCase| les classes.

\subsection{Les huit types primitifs}

\begin{tabz}{lXl}
  \headrow \thd{Type} \& \thd{Description} \& \thd{Exemple} \\
  \midrule
  \verb|string| \& S\'equence UTF-16 immuable \& \verb|"nour"|, \verb|'x'|, \verb|`a\${b}c`| \\
  \verb|number| \& Flottant 64 bits (~\$2^{53}\$ entiers exacts) \& \verb|42|, \verb|3.14|, \verb|1e3|, \verb|0xff|, \verb|1\_000\_000| \\
  \verb|bigint| \& Entier de taille arbitraire \& \verb|9007199254740993n| \\
  \verb|boolean| \& \verb|true| / \verb|false| \& \\
  \verb|undefined| \& Valeur non fournie \& \verb|let x;| \\
  \verb|null| \& Absence volontaire de valeur \& \verb|document.getElementById("x")| \\
  \verb|symbol| \& Cl\'e unique et immuable \& \verb|Symbol("id")| \\
\end{tabz}

Tout le reste est un \textbf{objet} (fonctions, tableaux, \verb|Map|, \verb|Date|, \verb|Error|\ldots).

\begin{codejs}
typeof 42;              // "number"
typeof "a";             // "string"
typeof true;            // "boolean"
typeof undefined;       // "undefined"
typeof Symbol();        // "symbol"
typeof 10n;             // "bigint"
typeof function () {};  // "function"
typeof {};              // "object"
typeof null;            // "object"  <-- BUG historique de 1995, jamais corrige. NE T'Y FIE PAS.
typeof [];              // "object"  <-- un tableau EST un objet
Array.isArray([]);      // true  <-- le VRAI test de tableau

Number.isInteger(5);        // true
Number.isNaN(NaN);          // true   (isNaN() global est faux : isNaN("toto") => true)
Number.isSafeInteger(2**53);// false
0.1 + 0.2 === 0.3;          // false  <-- arithmetique flottante : 0.30000000000000004
Math.abs(0.1 + 0.2 - 0.3) < Number.EPSILON;   // true : la bonne facon de comparer
\end{codejs}

\begin{piege}
\textbf{Trois pi\`eges de type, pour la vie.}
(1) \verb|typeof null === "object"|~: teste avec \verb|x === null|.
(2) \verb|0.1 + 0.2 !== 0.3|~: en mon\'etaire, travaille en centimes entiers ou avec une
biblioth\`eque d\'ecimale (\link{decimal.js}{https://github.com/MikeMcl/decimal.js/}).
(3) \verb|NaN| n'est \'egal \`a rien, pas m\^eme \`a lui-m\^eme~: utilise \verb|Number.isNaN(x)|.
\end{piege}

\section{Op\'erateurs}

\begin{codejs}
// Arithmetiques
7 / 2; 7 % 2; 2 ** 10; -(2); Math.trunc(7 / 2); Math.floor(-7 / 2);
(5.7).toFixed(2);              // "5.70"  -> RETOURNE UNE CHAINE
parseInt("42abc", 10); parseFloat("3.14"); Number("42"); +"42";

// Affectations combinees
let n = 10;
n += 5; n -= 2; n *= 3; n **= 2; n /= 4; n %= 7; n **= 2;
n++; ++n;                       // postfixe renvoie l'ancienne valeur, prefixe la nouvelle

// Comparaison --- la seule regle qui compte : TOUJOURS === et !==
5 == "5";        // true   (coercion : convertit avant de comparer)  INTERDIT
5 === "5";       // false  (comparaison stricte : type + valeur)
null == undefined;   // true
null === undefined;  // false
NaN === NaN;         // false

// Logiques
true && false;    // false  (court-circuit : le 2e operande n'est pas evalue)
true || false;    // true
!true;            // false
// Les formes modernes, tres utiles :
user?.profile?.avatar ?? "/default.png";   // chainage optionnel + alternative nulle
count ||= 10;     // assigne si falsy
config ??= {};    // assigne seulement si null/undefined  (le bon choix)

// Bit a bit (rare en web, courant en protocole/poids lourds)
5 & 3, 5 | 3, 5 ^ 3, ~5, 5 << 2, 5 >>> 1;

// Concatenation et template literal
"bonjour" + " " + nom;
`Bonjour ${nom}, tu as ${age} ans et ${[1, 2, 3].join(", ")} fruits.`;
\end{codejs}

\subsection{Conversion et \og ~falsy~ \fg{}}

Sept valeurs sont \textbf{falsy}~: \verb|false|, \verb|0|, \verb|-0|, \verb|0n|,
\verb|""|, \verb|null|, \verb|undefined|. Tout le reste est \textbf{truthy} -- y compris
\verb|[]|, \verb|{}| et \verb|"0"|.

\begin{codejs}
Boolean("");        // false
Boolean("0");       // true  <-- attention aux chaines non vides
if (tab.length) { /* ... */ }         // idiome courant
const nom = prenom || "Anonyme";      // remplace tout falsy (0, "", false)
const age = ageSaisi ?? 18;           // ne remplace que null/undefined  <-- RECOMMANDE
String(42); Number("42"); 42 + "";  +"42";
[1, "2", "3"] .map(Number);            // [1, 2, 3]
\end{codejs}

\section{Conditions et boucles}

\begin{codejs}
const note = 14;

if (note >= 16)       console.log("Tres bien");
else if (note >= 10)  console.log("Admis");
else                  console.log("Repasse");

// Expression conditionnelle (ternaire) : une VALEUR, pas un bloc
const mention = note >= 16 ? "tres bien" : note >= 10 ? "assez bien" : "insuffisant";

// switch : arrets obligatoires, sinon le code du cas suivant s'execute aussi
switch (jour) {
  case "samedi":
  case "dimanche":
    console.log("week-end");
    break;
  default:
    console.log("jour ouvrable");
}

// switch moderne : expression qui renvoie une valeur, correspondance par motifs
const tarif = switch (type) {
  case "enfant":        yield 5;
  case "adulte":        yield 12;
  default:              yield 8;
};                      // (actif a partir de ES2025 / Node 22+, sinon : ternaires ou objet)

const prix = { enfant: 5, adulte: 12, _default: 8 }[type] ?? 8;  // alternative simple et portative

for (let i = 0; i < 10; i++) { if (i === 3) continue; if (i === 8) break; }
let j = 0; while (j < 10) { j++; }
let k = 0; do { k++; } while (k < 10);   // execute au moins une fois

for (const v of [10, 20, 30]) {}          // itere sur les VALEURS (iterables)
for (const i in tab) {}                    // itere sur les CLES (chaines) -- a eviter sur tableaux
for (const [i, v] of tab.entries()) {}     // index + valeur
\end{codejs}

\begin{piege}
\verb|for...in| parcourt aussi les propri\'et\'es h\'erit\'ees \'enum\'erables et renvoie des
\emph{cha\^ines}. Sur un tableau, utilise \verb|for...of|, \verb|forEach| ou une boucle
indice. \verb|break| et \verb|continue| ne fonctionnent pas dans \verb|forEach|~: si tu as besoin
d'interrompre, c'est le signe qu'il te faut \verb|for...of|, \verb|some|, \verb|find| ou \verb|every|.
\end{piege}

\section{Fonctions}

\begin{codejs}
// Declaration (hoistee : utilisable avant sa definition)
function additionne(a, b = 0) { return a + b; }

// Expression de fonction (non hoistee)
const multiply = function (a, b) { return a * b; };

// Fleche : syntaxe compacte, pas de "this" propre
const square = (x) => x * x;
const toUser = (name, age) => ({ name, age, adult: age >= 18 });  // objet : parentheses obligatoires
const somme = (...args) => args.reduce((a, b) => a + b, 0);       // reste des arguments

// Arguments nommes + valeurs par defaut + valeurs de test
function envoi({
  destinataire,
  sujet = "(sans objet)",
  pieces = [],
  priorite = "normale",
} = {}) {
  if (!destinataire) throw new TypeError("destinataire requis");
  return { destinataire, sujet, pieces: pieces.length, priorite };
}
envoi({ destinataire: "nour@example.com", pieces: ["facture.pdf"] });

// Fonctions a rappels (callbacks) et fonctions d'ordre superieur
const appliquer = (fn, valeur) => fn(valeur);
const multiplieur = (facteur) => (x) => x * facteur;   // fabrique de fonctions
const triple = multiplieur(3); triple(5);               // 15

// R\'ecursivite : toujours prevoir un cas de base
function factorielle(n) {
  if (n < 0) throw new RangeError("n doit etre >= 0");
  return n <= 1 ? 1 : n * factorielle(n - 1);
}
\end{codejs}

\subsection{Port\'ee, hoisting, fermetures}

\begin{description}
  \item[Port\'ee de bloc] \verb|let| et \verb|const| vivent dans les accolades les plus proches.
        \verb|var| remonte jusqu'\`a la fonction~; c'est la source de bugs classiques dans les boucles.
  \item[Hoisting] Les d\'eclarations \verb|function| et \verb|var| sont \og ~remont\'ees~ \fg{} en m\'emoire.
        \verb|let|/\verb|const| existent mais sont dans la \og ~zone morte~ \fg{}~: y toucher avant
        l\`eve \verb|ReferenceError|.
  \item[Fermeture (closure)] Une fonction se souvient du environnement dans lequel elle a \'et\'e
        cr\'e\'ee. C'est le fondement de la vie priv\'ee en JS, des fabriques, de la m\'emo\"isation,
        des compteurs.
\end{description}

\begin{codejs}
// Le bug classique de var, a comprendre une fois pour toutes
for (var i = 0; i < 3; i++) setTimeout(() => console.log("var:", i)); // 3 3 3
for (let j = 0; j < 3; j++) setTimeout(() => console.log("let:", j)); // 0 1 2

// Fermeture : un compteur avec veritable encapsulation
function creeCompteur(demarrage = 0, pas = 1) {
  let valeur = demarrage;                       // prive : inaccessible de l'exterieur
  return {
    incrementer: () => (valeur += pas),
    lire: () => valeur,
    reinitialiser: () => (valeur = demarrage),
  };
}
const c = creeCompteur(10, 2);
c.incrementer(); c.incrementer(); c.lire();        // 14
c.valeur;                                          // undefined -- c'est tout l'interet

// Memoisation grace a une fermeture
function memoise(fn, cle = (...a) => JSON.stringify(a)) {
  const cache = new Map();
  return (...args) => {
    const k = cle(...args);
    if (!cache.has(k)) cache.set(k, fn(...args));
    return cache.get(k);
  };
}
const fib = memoise(function f(n) { return n < 2 ? n : f(n - 1) + f(n - 2); });
fib(70);   // immediat
\end{codejs}

\section{Le module commun : erreurs, piles, d\'ebogage}

\begin{codejs}
function diviser(a, b) {
  if (b === 0) throw new RangeError("division par zero");
  return a / b;
}
try {
  const r = diviser(10, 0);
} catch (err) {
  if (err instanceof RangeError) console.error("Domaine invalide :", err.message);
  else throw err;                        // ne JAMAIS avaler une erreur inconnue
} finally {
  console.log("toujours execute (fermer une connexion, logger, etc.)");
}

// Empilement des causes (ES2022) : precieux en production
try {
  JSON.parse("{oops");
} catch (cause) {
  throw new Error("Configuration invalide", { cause });
}
\end{codejs}

D\'eboguer proprement~: un point d'arr\^et dans l'onglet \clic{Sources} de Chrome (\keys{F12}),
puis \keys{F10} (pas \`a pas principal), \keys{F11} (entrer), \keys{F8} (continuer)~;
\verb|debugger;| dans le code ; \verb|console.trace()| ; \verb|console.log(new Error().stack)|.
La quasi-totalit\'e des d\'ebutants abuse de \verb|console.log|~: trois points d'arr\^et bien plac\'es
valent cent traces.

\begin{exercice}
\begin{enumerate}
  \item \textbf{FizzBuzz}~: affiche de 1 \`a 100, \og ~Fizz~ \fg{} si multiple de 3, \og ~Buzz~ \fg{} de 5,
        \og ~FizzBuzz~ \fg{} des deux. \'Ecris-le trois fois~: avec \verb|if|, avec un ternaire imbriqu\'e,
        avec une composition de fonctions.
  \item \textbf{Pair ou impair}, \textbf{plus grand de trois nombres}, \textbf{reverse d'une cha\^ine}
        (avec et sans \verb|split/reverse/join|).
  \item \textbf{Palindrome normalis\'e}~: \verb|"\'Esope reste ici et se repose"| renvoie \verb|true|.
        Retirer ponctuation, espaces, accents (\verb|String.prototype.normalize("NFD")|).
  \item \textbf{Nombre premier}~: fonction \verb|estPremier(n)| avec boucle jusqu'\`a
        \$\sqrt{n}\$~; puis \verb|premiersInfrieursA(100)|.
  \item \textbf{Temperature}~: \verb|celsiusVersFahrenheit|, \verb|fahrenheitVersCelsius|,
        et le sens inverse, \'ecrites comme fleches. V\'erifie que aller-retour redonne la valeur.
  \item \textbf{Compteur de clics}~: avec la fermeture \verb|creeCompteur|, simule 5 clics
        et affiche la valeur \`a chaque fois.
  \item \textbf{Devine le r\'esultat} avant d'ex\'ecuter~: \verb|1 + "2"|, \verb|"5" - 3|,
        \verb|"a" + + "b"|, \verb|[] + []|, \verb|[] + {}|, \verb|null + 1|,
        \verb|undefined + 1|, \verb|true + true|, \verb|2 + 3 + "a"|, \verb|"a" + 2 + 3|.
\end{enumerate}
\end{exercice}

\begin{ressources}
\begin{itemize}
  \item MDN --- Guide JavaScript (la bible, en fran\c{c}ais) --- \url{https://developer.mozilla.org/fr/docs/Web/JavaScript/Guide}
  \item \emph{JavaScript.info}, le meilleur cours libre du web --- \url{https://javascript.info/}
  \item \emph{Eloquent JavaScript}, 3\textsuperscript{e} \'edition --- \url{https://eloquentjavascript.net/}
  \item Table de compatibilit\'e des fonctionnalit\'es --- \url{https://compat-table.github.io/compat-table/es2016plus/}
  \item \emph{You Don't Know JS Yet} (Deepankar) --- \url{https://github.com/getify/You-Dont-Know-JS}
\end{itemize}
\end{ressources}

\index{JavaScript!bases}\index{variables}\index{closure}
